\documentclass[10pt,a4paper]{amsart}
\usepackage{amsmath,amssymb,amsthm,graphicx}
\usepackage{verbatim}
\newenvironment{reference}{\comment}{\endcomment}
%\newenvironment{reference}{}{}
\newenvironment{slogan}{\comment}{\endcomment}
\newenvironment{history}{\comment}{\endcomment}

% For commutative diagrams we use Xy-pic
\usepackage[all]{xy}

% We use 2cell for 2-commutative diagrams.
\xyoption{2cell}
% Theorem environments.
%
\theoremstyle{plain}
\newtheorem{theorem}[subsection]{Theorem}
\newtheorem{proposition}[subsection]{Proposition}
\newtheorem{lemma}[subsection]{Lemma}

\theoremstyle{definition}
\newtheorem{definition}[subsection]{Definition}
\newtheorem{example}[subsection]{Example}
\newtheorem{exercise}[subsection]{Exercise}
\newtheorem{situation}[subsection]{Situation}

\theoremstyle{remark}
\newtheorem{remark}[subsection]{Remark}
\newtheorem{remarks}[subsection]{Remarks}

\numberwithin{equation}{subsection}

% Macros
%
\def\lim{\mathop{\mathrm{lim}}\nolimits}
\def\colim{\mathop{\mathrm{colim}}\nolimits}
\def\Spec{\mathop{\mathrm{Spec}}}
\def\Hom{\mathop{\mathrm{Hom}}\nolimits}
\def\Ext{\mathop{\mathrm{Ext}}\nolimits}
\def\SheafHom{\mathop{\mathcal{H}\!\mathit{om}}\nolimits}
\def\SheafExt{\mathop{\mathcal{E}\!\mathit{xt}}\nolimits}
\def\Sch{\mathit{Sch}}
\def\Mor{\mathop{\mathrm{Mor}}\nolimits}
\def\Ob{\mathop{\mathrm{Ob}}\nolimits}
\def\Sh{\mathop{\mathit{Sh}}\nolimits}
\def\NL{\mathop{N\!L}\nolimits}
\def\CH{\mathop{\mathrm{CH}}\nolimits}
\def\proetale{{pro\text{-}\acute{e}tale}}
\def\etale{{\acute{e}tale}}
\def\QCoh{\mathit{QCoh}}
\def\Ker{\mathop{\mathrm{Ker}}}
\def\Im{\mathop{\mathrm{Im}}}
\def\Coker{\mathop{\mathrm{Coker}}}
\def\Coim{\mathop{\mathrm{Coim}}}

% Boxtimes
%
\DeclareMathSymbol{\boxtimes}{\mathbin}{AMSa}{"02}

%
% Macros for moduli stacks/spaces
%
\def\QCohstack{\mathcal{QC}\!\mathit{oh}}
\def\Cohstack{\mathcal{C}\!\mathit{oh}}
\def\Spacesstack{\mathcal{S}\!\mathit{paces}}
\def\Quotfunctor{\mathrm{Quot}}
\def\Hilbfunctor{\mathrm{Hilb}}
\def\Curvesstack{\mathcal{C}\!\mathit{urves}}
\def\Polarizedstack{\mathcal{P}\!\mathit{olarized}}
\def\Complexesstack{\mathcal{C}\!\mathit{omplexes}}
% \Pic is the operator that assigns to X its picard group, usage \Pic(X)
% \Picardstack_{X/B} denotes the Picard stack of X over B
% \Picardfunctor_{X/B} denotes the Picard functor of X over B
\def\Pic{\mathop{\mathrm{Pic}}\nolimits}
\def\Picardstack{\mathcal{P}\!\mathit{ic}}
\def\Picardfunctor{\mathrm{Pic}}
\def\Deformationcategory{\mathcal{D}\!\mathit{ef}}
\usepackage{cleveref}
\crefname{originallemma}{original lemma}{original lemmas}
\Crefname{originallemma}{Original lemma}{Original lemmas}
\makeatletter
\newlabel{derived-lemma-compact-objects-subcategory}{{09QH}{source}{}{sourcepointer-1}{}}
\newlabel{derived-lemma-compact-objects-subcategory@cref}{{[originallemma][0][]09QH}{[1][0][]source}{}{}{}}
\makeatother
\crefname{originallemma}{original lemma}{original lemmas}
\Crefname{originallemma}{Original lemma}{Original lemmas}
\makeatletter
\newlabel{derived-lemma-composition-zero}{{0146}{source}{}{sourcepointer-2}{}}
\newlabel{derived-lemma-composition-zero@cref}{{[originallemma][0][]0146}{[1][0][]source}{}{}{}}
\makeatother
\crefname{originallemma}{original lemma}{original lemmas}
\Crefname{originallemma}{Original lemma}{Original lemmas}
\makeatletter
\newlabel{derived-lemma-split}{{05QT}{source}{}{sourcepointer-3}{}}
\newlabel{derived-lemma-split@cref}{{[originallemma][0][]05QT}{[1][0][]source}{}{}{}}
\makeatother
\crefname{originalreference}{original reference}{original references}
\Crefname{originalreference}{Original reference}{Original references}
\makeatletter
\newlabel{lemma-derived-products}{{07D9}{source}{}{sourcepointer-4}{}}
\newlabel{lemma-derived-products@cref}{{[originalreference][0][]07D9}{[1][0][]source}{}{}{}}
\makeatother
\crefname{originalreference}{original reference}{original references}
\Crefname{originalreference}{Original reference}{Original references}
\makeatletter
\newlabel{lemma-factor-through-nicer}{{09R2}{source}{}{sourcepointer-5}{}}
\newlabel{lemma-factor-through-nicer@cref}{{[originalreference][0][]09R2}{[1][0][]source}{}{}{}}
\makeatother
\crefname{originalreference}{original reference}{original references}
\Crefname{originalreference}{Original reference}{Original references}
\makeatletter
\newlabel{lemma-hom-derived}{{09KY}{source}{}{sourcepointer-6}{}}
\newlabel{lemma-hom-derived@cref}{{[originalreference][0][]09KY}{[1][0][]source}{}{}{}}
\makeatother
\crefname{originalreference}{original reference}{original references}
\Crefname{originalreference}{Original reference}{Original references}
\makeatletter
\newlabel{lemma-property-P-sequence}{{09KL}{source}{}{sourcepointer-7}{}}
\newlabel{lemma-property-P-sequence@cref}{{[originalreference][0][]09KL}{[1][0][]source}{}{}{}}
\makeatother
\crefname{originalreference}{original reference}{original references}
\Crefname{originalreference}{Original reference}{Original references}
\makeatletter
\newlabel{lemma-resolve}{{09KP}{source}{}{sourcepointer-8}{}}
\newlabel{lemma-resolve@cref}{{[originalreference][0][]09KP}{[1][0][]source}{}{}{}}
\makeatother
\title{Compact objects of the derived category of a differential graded algebra are the direct summands of modules with a finite filtration by shifts of the algebra}
\author{The Stacks Project Authors}
\date{stacks-project:dga.tex:local-snapshot-20261012; GNU FDL 1.2 or later}
\begin{document}
\maketitle
\noindent\textit{Editorial note.} All delivered mathematical text is the author's own, copied from the pinned Stacks Project chapter file sources/stacks-dga/pinned-original-source.tex; no statement or proof step was rewritten, reordered or removed.\par\medskip
\noindent\textit{Editorial note.} The counted proof is the author's complete argument: the property-(P) representation, the filtration triangle with its boundary analysis via the difference map, the two compactness factorization steps, the upgrade through lemma-factor-through-nicer, the converse reduction through property (P) and the cohomologicality of Hom, and the final shift case are all retained verbatim (the author marks the shift case as clear; Hom(A[k], -) being a cohomology functor is the standard one-step check).\par\medskip
\noindent\textit{Editorial note.} The label proposition-compact names the counted statement, and the proof environment immediately following it in the source is delivered verbatim as its proof.\par\medskip
\noindent\textit{Editorial note.} The counted statement is carried by the author's own proposition environment in the chapter source.\par\medskip
\noindent\textit{Allowed standard prerequisites (editorial).}\begin{itemize}
\item Property (P) resolutions and the associated triangle. Same chapter, Lemmas lemma-resolve, lemma-property-P-sequence and lemma-factor-through-nicer. Lets a compact object be represented by a module P with property (P), provides the distinguished triangle of the filtration whose boundary map the compactness argument analyzes, and upgrades a factorization through a finite partial sum to the filtration-form of part (2).
\item Compactness formalism in derived categories. Derived Categories, Lemmas derived-lemma-composition-zero, derived-lemma-split and derived-lemma-compact-objects-subcategory; Same chapter, Lemma lemma-hom-derived and lemma-derived-products. Supplies the vanishing of the composed boundary map, the splitting of the triangle once the boundary vanishes, the reduction to direct summands, the identification of derived Hom with homotopy Hom for property-(P) modules, and the description of direct sums in the derived category.
\end{itemize}\medskip
\ifdefined\hypertarget\hypertarget{sourcepointer-1}{}\fi
\noindent\textit{Source reference.} Original lemma 09QH: sources/stacks-chow/tags.txt, source lines 12285--12285. Prerequisite pointer invoked at proof line(s) 6546; it names a standard fact of a sibling chapter that the counted argument uses, and every step of the counted proof built on it is retained verbatim in the delivered body.\par\medskip
\ifdefined\hypertarget\hypertarget{sourcepointer-2}{}\fi
\noindent\textit{Source reference.} Original lemma 0146: sources/stacks-chow/tags.txt, source lines 1461--1461. Prerequisite pointer invoked at proof line(s) 6528; it names a standard fact of a sibling chapter that the counted argument uses, and every step of the counted proof built on it is retained verbatim in the delivered body.\par\medskip
\ifdefined\hypertarget\hypertarget{sourcepointer-3}{}\fi
\noindent\textit{Source reference.} Original lemma 05QT: sources/stacks-chow/tags.txt, source lines 7387--7387. Prerequisite pointer invoked at proof line(s) 6534; it names a standard fact of a sibling chapter that the counted argument uses, and every step of the counted proof built on it is retained verbatim in the delivered body.\par\medskip
\ifdefined\hypertarget\hypertarget{sourcepointer-4}{}\fi
\noindent\textit{Source reference.} Original reference 07D9: sources/stacks-dga/pinned-original-source.tex, source lines 2895--2895. Prerequisite pointer invoked at proof line(s) 6554; it names a standard fact of the same chapter that the counted argument uses, and every step of the counted proof built on it is retained verbatim in the delivered body.\par\medskip
\ifdefined\hypertarget\hypertarget{sourcepointer-5}{}\fi
\noindent\textit{Source reference.} Original reference 09R2: sources/stacks-dga/pinned-original-source.tex, source lines 6434--6434. Prerequisite pointer invoked at proof line(s) 6539; it names a standard fact of the same chapter that the counted argument uses, and every step of the counted proof built on it is retained verbatim in the delivered body.\par\medskip
\ifdefined\hypertarget\hypertarget{sourcepointer-6}{}\fi
\noindent\textit{Source reference.} Original reference 09KY: sources/stacks-dga/pinned-original-source.tex, source lines 2850--2850. Prerequisite pointer invoked at proof line(s) 6552; it names a standard fact of the same chapter that the counted argument uses, and every step of the counted proof built on it is retained verbatim in the delivered body.\par\medskip
\ifdefined\hypertarget\hypertarget{sourcepointer-7}{}\fi
\noindent\textit{Source reference.} Original reference 09KL: sources/stacks-dga/pinned-original-source.tex, source lines 2425--2425. Prerequisite pointer invoked at proof line(s) 6524; it names a standard fact of the same chapter that the counted argument uses, and every step of the counted proof built on it is retained verbatim in the delivered body.\par\medskip
\ifdefined\hypertarget\hypertarget{sourcepointer-8}{}\fi
\noindent\textit{Source reference.} Original reference 09KP: sources/stacks-dga/pinned-original-source.tex, source lines 2519--2519. Prerequisite pointer invoked at proof line(s) 6516; it names a standard fact of the same chapter that the counted argument uses, and every step of the counted proof built on it is retained verbatim in the delivered body.\par\medskip
\begin{proposition}
\label{proposition-compact}
Let $(A, \text{d})$ be a differential graded algebra. Let $E$ be an
object of $D(A, \text{d})$. Then the following are equivalent
\begin{enumerate}
\item $E$ is a compact object,
\item $E$ is a direct summand of an object of $D(A, \text{d})$
which is represented by a differential graded module $P$ which
has a finite filtration $F_\bullet$ by differential graded submodules
such that $F_iP/F_{i - 1}P$ are finite direct sums of shifts of $A$.
\end{enumerate}
\end{proposition}

\begin{proof}
Assume $E$ is compact. By Lemma \ref{lemma-resolve} we may assume that $E$
is represented by a differential graded $A$-module $P$ with property (P).
Consider the distinguished triangle
$$
\bigoplus F_iP \to \bigoplus F_iP \to P
\xrightarrow{\delta} \bigoplus F_iP[1]
$$
coming from the admissible short exact sequence of
Lemma \ref{lemma-property-P-sequence}. Since $E$ is compact we have
$\delta = \sum_{i = 1, \ldots, n} \delta_i$ for some
$\delta_i : P \to F_iP[1]$. Since the composition of $\delta$
with the map $\bigoplus F_iP[1] \to \bigoplus F_iP[1]$ is zero
(Derived Categories, Lemma \ref{derived-lemma-composition-zero})
it follows that $\delta = 0$ (follows as $\bigoplus F_iP \to \bigoplus F_iP$
maps the summand $F_iP$ via the difference of $\text{id}$ and the inclusion
map into $F_{i - 1}P$).
Thus we see that the identity on $E$ factors through
$\bigoplus F_iP$ in $D(A, \text{d})$ (by
Derived Categories, Lemma \ref{derived-lemma-split}).
Next, we use that $P$ is compact again to see that the map
$E \to \bigoplus F_iP$ factors through $\bigoplus_{i = 1, \ldots, n} F_iP$
for some $n$. In other words, the identity on $E$ factors through
$\bigoplus_{i = 1, \ldots, n} F_iP$. By
Lemma \ref{lemma-factor-through-nicer}
we see that the identity of $E$ factors as $E \to P \to E$
where $P$ is as in part (2) of the statement of the lemma.
In other words, we have proven that (1) implies (2).

\medskip\noindent
Assume (2). By
Derived Categories, Lemma \ref{derived-lemma-compact-objects-subcategory}
it suffices to show that $P$ gives a compact object. Observe that
$P$ has property (P), hence we have
$$
\Hom_{D(A, \text{d})}(P, M) = \Hom_{K(A, \text{d})}(P, M)
$$
for any differential graded module $M$ by Lemma \ref{lemma-hom-derived}.
As direct sums in $D(A, \text{d})$ are given by direct sums of
graded modules (Lemma \ref{lemma-derived-products}) we reduce
to showing that $\Hom_{K(A, \text{d})}(P, M)$ commutes with direct
sums. Using that $K(A, \text{d})$ is a triangulated category,
that $\Hom$ is a cohomological functor in the first
variable, and the filtration on $P$, we reduce to the case that
$P$ is a finite direct sum of shifts of $A$. Thus we reduce to
the case $P = A[k]$ which is clear.
\end{proof}

\end{document}
